\documentclass[12pt]{article}
\usepackage[slovene]{babel}
\usepackage{lmodern}

\usepackage{minted}
\setminted{fontsize=\small, linenos, }
% % tcolorbox / tcblisting
% % 
\usepackage{booktabs,array}
\usepackage{multirow}

\author{France Prešeren}
\title{Zdravljica}
\date{\today}  

\begin{document}
  \maketitle
  To je moja prva datoteka v LaTeXu. V tem bom napisal ljubezensko pismo Juliji.
  
  \vspace{2cm}

  "` "' Več kot 50 \% površine Slovenije pokriva gozd "> "< .

  \glqq \grqq

  \vspace{2cm}


  Več kot 90 \% Slovencev zasluži več kot 500 € na mesec, kar je približno 600 \$.

  \vspace{2cm}

  Ni pomembno, če za besedo
  stoji en ali pa več           presledkov.
  Prelom vrstice brez praznih 
        se obravnava kot presledek.


  Prazna vrstica začenja nov odstavek. 
  %
  Prazen komentar ni prazna vrstica.

\clearpage

\noindent Živé naj vsi naródi,\\
ki hrepené dočakat' dan,\newline
da koder sonce hodi,\\
prepir iz svéta bo pregnan,\newline
da rojak\\
prost bo vsak,\\
ne vrag, le sosed bo mejak!

\vspace{2cm}
Vpliva samo na \textit{argument}.
Po \texttt{\textbackslash itshape} \itshape
je vse v poševnem tisku,
dokler ne sledi  \upshape \texttt{\textbackslash upshape}.
Zdaj je vse tako, kot je bilo prej.



\vspace{2cm}

\begin{center}
 \begin{tabular}{lll}
\toprule
Pisava & Ukaz (za argument) & Ukaz (switch) \\
\midrule
\textrm{pokončna serifna pisava} & \mintinline{latex}{\textrm{<besedilo>}} & \mintinline{latex}{\rmfamily} \\
\textsf{gladka pisava} & \mintinline{latex}{\textsf{<besedilo>}} & \mintinline{latex}{\sffamily} \\
\texttt{strojepisna pisava} & \mintinline{latex}{\texttt{<besedilo>}} & \mintinline{latex}{\ttfamily} \\ \midrule
\textmd{srednja debelina} & \mintinline{latex}{\textmd{<besedilo>}} & \mintinline{latex}{\mdseries} \\
\textbf{krepka pisava} & \mintinline{latex}{\textbf{<besedilo>}} & \mintinline{latex}{\bfseries} \\ \midrule
\textit{poševna pisava} & \mintinline{latex}{\textit{<besedilo>}} & \mintinline{latex}{\itshape} \\
\textsl{ležeča pisava} & \mintinline{latex}{\textsl{<besedilo>}} & \mintinline{latex}{\slshape} \\
\textsc{male velike črke} & \mintinline{latex}{\textsc{<besedilo>}} & \mintinline{latex}{\scshape} \\\midrule
\textup{pokončna} & \mintinline{latex}{\textup{<besedilo>}} & \mintinline{latex}{\upshape} \\
\textup{velika-mala} & \mintinline{latex}{\textulc{<besedilo>}} & \mintinline{latex}{\ulcshape} \\
\textnormal{privzeta pisava} & \mintinline{latex}{\textnormal{<besedilo>}} & \mintinline{latex}{\normalfont} \\
\bottomrule
\end{tabular}
\end{center}

\scshape
does \ulcshape this work? 


\clearpage

% \vspace{2cm}

To je \underline{podčrtano} besedilo. 
Izgleda \underline{malo} \underline{grdo}, 
zato ga ne uporabljajmo pogosto.

To je \emph{poudarjeno} besedilo.

\textit{To je poševno besedilo, v 
katerem je \emph{poudarjeno} besedilo.}

\emph{Ta \emph{poudarjeno besedilo} 
je obdan s poudarjenim besedilom \dots 
ali je \emph{res} poudarjen?}


\vspace{2cm}

Praviloma pisave izbiramo premišljeno, ali pa izbor kar prepustimo LaTeXu.
Besedilo lahko \emph{poudarimo} ali zapišemo \textbf{krepko}, ali pa \emph{\textbf{oboje skupaj}}.
Kaj se zgodi, če \emph{uporabimo poudarjeno \emph{znotraj} poudarjenega besedila}?
Seveda \textbf{lahko nastavimo tudi \textnormal{običajno} pisavo}.
Besedilo lahko tudi \underline{podčrtamo}, vendar tega ne priporočamo, \underline{ker} \underline{je grdo}.

Pišemo lahko tudi v \textsf{sans-serifni pisavi} ali pa \textsc{z malimi velikimi črkami},
denimo \textsc{SageMath}. 
\textsl{Ležeča pisava} {ni ista reč kot} \textit{poševna pisava}.
Poševna pisava uporablja \textit{drugačne glife}, ležeča pisava pa uporablja \textsl{iste glife}, vendar ležečo. 
Razlika je res očitna, če pogledate \textit{poševno pisano črko a} poleg \textsl{ležeče črke a}.

Včasih uporabimo tudi \texttt{pisavo fiksne širine}, v kateri so vsi znaki enako široki.


\vspace{2cm}

\begin{tabular}{p{0.4\textwidth}p{0.55\textwidth}}
        % \toprule
        Ukaz  & Velikost                               \\
        \hline
        \mintinline{latex}{{\tiny}}  & {\tiny drobna pisava} \\                      
        \mintinline{latex}{{\scriptsize}} & {\scriptsize velikost indeksov}           \\
        \mintinline{latex}{{\footnotesize}}           & {\footnotesize  velikost opomb pod črto}       \\
        \mintinline{latex}{{\small}}    & {\small majhna pisava}                     \\
        \mintinline{latex}{{\normalsize}}             & {\normalsize  normalna velikost}              \\
        \mintinline{latex}{{\large}}                  & {\large nekoliko večji znaki}                     \\
        \mintinline{latex}{{\Large}}  & {\Large veliki znaki}                      \\
        \mintinline{latex}{{\LARGE}}                  &{\LARGE zelo veliki znaki}                  \\
        \multirow{2}{*}{\mintinline{latex}{{\huge}} } & \multirow{2}{*}{{\huge ogromni znaki}}        \\
        & \\
        \multirow{2}{*}{\mintinline{latex}{{\Huge}} } & \multirow{2.2}{*}{{\Huge največji znaki} }   \\
        & \\
      \end{tabular}

\vspace{2cm}
      \begin{quote}
        \underline{\textbf{Pom\textsc{nite}\Huge!}} \textit{Čim}
        \textsf{V\textbf{\LARGE E}}\textit{č} pisav \Huge uporabljate
         \footnotesize \textbf{v} vašem \small \texttt{dokumentu},
        \large \textit{tem} \normalsize lažje \textsc{berljiv} in
        \textsl{\textsf{lepši} pos\large t\Large a\LARGE n\huge e}.
      \end{quote}

\vspace{2cm}

V slovenščini vezaj (-) brez presledkov povezuje sestavine besed (slovensko-angleški slovar, C-vitamin), medtem ko pomišljaj (--) s presledki razmejuje ali poudarja miselne enote (To je bilo -- verjeli ali ne -- res.) ter označuje obseg (strani 15--20).
V slovenščini uporabljamo več vrst narekovajev. 
Najpogostejši so dvojni spodnji in zgornji ("` "'), npr. "`To je primer."' 
Uporabljajo se v knjižnem jeziku in uradnih besedilih. 
Pogosti so tudi dvojni zgornji (`` '' ), npr. ``Tako piše v članku.'', zlasti v elektronskih besedilih. Včasih naletimo na špičaste narekovaje ("> "<), npr. ">Tako so govorili."<, ki imajo dolgo tradicijo v slovenski tipografiji. 
Vsi trije tipi so pravilni, pomembno pa je, da v enem besedilu uporabljamo enotno obliko.

Ko pišemo besedilo v \LaTeX u, moramo paziti, da uporabljamo pravilne simbole. Poleg zgoraj navedenih simbolov ima \LaTeX tudi poseben ukaz za \dots, in sicer \dots{}  (prim. ...).

\vspace{2cm}

Ukazi, kot je \texttt{\textit{argument}},
  vplivajo samo na \textit{argument}. 
  Če uporabimo različico switch, se vse, 
  kar je po ukazom \texttt{\itshape} \itshape,  
  spremeni in ostane tako, dokler se ne pojavi drug \texttt{upshape}. 
  Zdaj je vse tako, kot je bilo prej.

\end{document}